% =============================================================
%  ch05_kinematik.tex  --  Kinematik, Dynamik und Trajektorien
% =============================================================

\chapter{Kinematik, Dynamik und Trajektorien}

\section{Vorwärtskinematik}
Aus den Gelenkwinkeln $\bm{q}$ die Endeffektor-Pose berechnen. Klassisch über
die Kette homogener $4\times4$-Transformationen (Denavit-Hartenberg-Parameter).
Der Rotationsblock jeder Transformation ist eine $3\times3$-Matrix; intern
äquivalent zur Quaternion-Darstellung.

\section{Inverse Kinematik}
Aus der gewünschten Pose die Gelenkwinkel finden. Numerisch über die
\textbf{Jacobi-Matrix} $\bm{J}$, die Gelenkgeschwindigkeiten auf
Endeffektorgeschwindigkeiten abbildet ($\dot{\bm{x}} = \bm{J}\dot{\bm{q}}$).
Iterativ:
\begin{equation}
\Delta\bm{q} = \bm{J}^{\dagger}\,\bm{e},
\qquad
\bm{e} = \begin{pmatrix}\bm{e}_{\text{pos}} \\ \bm{e}_{\text{ori}}\end{pmatrix}
\end{equation}
mit Pseudo-Inverser $\bm{J}^\dagger$ (oder Damped-Least-Squares nahe
Singularitäten). Der \textbf{Orientierungsfehler} $\bm{e}_{\text{ori}}$ wird als
Vektorteil des Fehlerquaternions $q_{\text{soll}}\otimes q_{\text{ist}}^{-1}$
gebildet -- Euler-Differenzen sind hier wegen Gimbal Lock unbrauchbar. In
ROS\,2 löst \textbf{MoveIt\,2} (mit KDL/TracIK) das für dich.

\section{Dynamik}
Das Bewegungsgesetz eines Roboterarms:
\begin{equation}
\bm{M}(\bm{q})\ddot{\bm{q}} + \bm{C}(\bm{q},\dot{\bm{q}})\dot{\bm{q}}
+ \bm{g}(\bm{q}) = \bm{\tau}
\end{equation}
$\bm{M}$ = Massenmatrix, $\bm{C}$ = Coriolis/Zentrifugal,
$\bm{g}$ = Gravitation, $\bm{\tau}$ = Gelenkmomente. Für die Simulation ist
entscheidend, dass du jedem Link in der URDF korrekte \textbf{Masse und
Trägheitstensoren} gibst -- sonst rechnet die Physik-Engine Unsinn.

\section{Trajektorienplanung}
\begin{itemize}
\item \textbf{Gelenkraum:} Kubische/quintische Polynome oder ein
  \textbf{trapezförmiges Geschwindigkeitsprofil} (beschleunigen -- konstant --
  abbremsen) erzeugen ruckarme, ableitungsstetige Verläufe.
\item \textbf{Kartesischer Raum:} Position linear interpolieren, Orientierung
  per \textbf{SLERP} (gleichmäßige Winkelgeschwindigkeit). Für glatte
  Mehrpunktbahnen \textbf{SQUAD}.
\end{itemize}
In ROS\,2 transportiert \code{trajectory\_msgs/JointTrajectory} eine Bahn
(Punkte mit Position/Geschwindigkeit/Zeit), und der
\code{joint\_trajectory\_controller} fährt sie ab.
